% part: intuitionistic-logic
% chapter: soundness-completeness
% section: truth-lemma

\documentclass[../../../include/open-logic-section]{subfiles}

\begin{document}

\olfileid{int}{sc}{tru}

\olsection{لمّة الصدق}

تربط اللمّة الآتية بين الإشباع في النموذج القانوني وبين ال!!{formula}s
التي هي !!{element}s في المجموعة الأولية~$\Delta$.

\begin{lem}\ollabel{lem:truth}
  إذا كانت~$\Delta$ أولية، كان~$\mSat{M(\Delta)}{!A}[\sigma]$ إذا
  وفقط إذا كان~$\Delta(\sigma) \Proves !A$.
\end{lem}

\begin{proof}
  بالاستقراء في~$!A$.
  \begin{enumerate}
    \item \indcase{!A}{\lfalse}{بما أن~$\Delta(\sigma)$ أولية، فهي
      متسقة، ولذلك~$\Delta(\sigma) \Proves/ \indfrm$. وبحسب التعريف،
      $\mSat/{M(\Delta)}{\indfrm}[\sigma]$.}
    \item \indcase{!A}{p}{بحسب تعريف~$\mSat{{}}{}$،
      $\mSat{M(\Delta)}{\indfrm}[\sigma]$ إذا وفقط إذا كان
      $\sigma \in V(p)$؛ أي إذا وفقط إذا كان
      $\Delta(\sigma) \Proves \indfrm$.}
    \item \indcase{!A}{\lnot !B}{هذه هي حالة الشرط مع نتيجة
      $\lfalse$، ويمكن تفصيلها على النحو الآتي. افترض أولًا
      $\Delta(\sigma)\Proves\lnot !B$. ولكل~$\sigma'$ بحيث
      $R\sigma\sigma'$، لدينا
      $\Delta(\sigma)\subseteq\Delta(\sigma')$، ومن ثم
      $\Delta(\sigma')\Proves\lnot !B$. ولو كان
      $\mSat{M(\Delta)}{!B}[\sigma']$، لأعطى فرض الاستقراء
      $\Delta(\sigma')\Proves !B$، فيناقض ذلك اتساق
      $\Delta(\sigma')$. ولذلك لا يصدق~$!B$ عند أي امتداد~$\sigma'$،
      أي~$\mSat{M(\Delta)}{\lnot !B}[\sigma]$.

      وللاتجاه العكسي نبرهن النقيض المقابل. افترض
      $\Delta(\sigma)\Proves/\lnot !B$. عندئذ
      $\Delta(\sigma)\cup\{!B\}\Proves/\lfalse$؛ إذ لو أمكن اشتقاق
      $\lfalse$ من إضافة~$!B$، لأمكن إدخال النفي واشتقاق
      $\lnot !B$ من~$\Delta(\sigma)$. ولتكن
      $\tuple{!B,\lfalse}=\tuple{!B_n,!C_n}$ لبعض~$n$. عندئذ
      $\Delta(\sigma.n)=(\Delta(\sigma)\cup\{!B\})^*$، ولذلك
      $\Delta(\sigma.n)\Proves !B$. وبفرض الاستقراء،
      $\mSat{M(\Delta)}{!B}[\sigma.n]$. ولأن
      $R\sigma(\sigma.n)$، يكون
      $\mSat/{M(\Delta)}{\lnot !B}[\sigma]$.}
    \item \indcase{!A}{!B \land
      !C}{$\mSat{M(\Delta)}{\indfrm}[\sigma]$ إذا وفقط إذا كان
      $\mSat{M(\Delta)}{!B}[\sigma]$ و
      $\mSat{M(\Delta)}{!C}[\sigma]$. وبفرض الاستقراء،
      $\mSat{M(\Delta)}{!B}[\sigma]$ إذا وفقط إذا كان
      $\Delta(\sigma) \Proves !B$، وبالمثل لـ~$!C$. لكن
      $\Delta(\sigma) \Proves !B$ و$\Delta(\sigma) \Proves !C$ إذا
      وفقط إذا كان~$\Delta(\sigma) \Proves \indfrm$.}
    \item \indcase{!A}{!B \lor
      !C}{$\mSat{M(\Delta)}{\indfrm}[\sigma]$ إذا وفقط إذا كان
      $\mSat{M(\Delta)}{!B}[\sigma]$ أو
      $\mSat{M(\Delta)}{!C}[\sigma]$. وبفرض الاستقراء، يتحقق ذلك إذا
      وفقط إذا كان~$\Delta(\sigma) \Proves !B$ أو
      $\Delta(\sigma) \Proves !C$. وعلينا أن نثبت أن هذا بدوره يتحقق
      إذا وفقط إذا كان~$\Delta(\sigma) \Proves \indfrm$. والاتجاه من
      اليسار إلى اليمين واضح، أما الاتجاه من اليمين إلى اليسار فينتج
      من كون~$\Delta(\sigma)$ أولية.}
    \item \indcase{!A}{!B \lif !C}{نبرهن أولًا النقيض المقابل للاتجاه
      من اليسار إلى اليمين. افترض
      $\Delta(\sigma) \Proves/ !B \lif !C$. عندئذ أيضًا
      $\Delta(\sigma) \cup \{!B\} \Proves/ !C$. ولأن
      $\tuple{!B, !C}$ هي~$\tuple{!B_n, !C_n}$ لبعض~$n$، يكون
      $\Delta(\sigma.n) = (\Delta(\sigma) \cup \{!B\})^*$، ويكون
      $\Delta(\sigma.n) \Proves !B$ لكن
      $\Delta(\sigma.n) \Proves/ !C$. وبفرض الاستقراء،
      $\mSat{M(\Delta)}{!B}[\sigma.n]$ و
      $\mSat/{M(\Delta)}{!C}[\sigma.n]$. ولأن
      $R\sigma(\sigma.n)$، فهذا يعني
      $\mSat/{M(\Delta)}{\indfrm}[\sigma]$.

      افترض الآن~$\Delta(\sigma) \Proves !B \lif !C$، ولتكن
      $R\sigma\sigma'$. وبما أن
      $\Delta(\sigma) \subseteq \Delta(\sigma')$، يكون: إذا كان
      $\Delta(\sigma') \Proves !B$، كان
      $\Delta(\sigma') \Proves !C$. وبعبارة أخرى، لكل~$\sigma'$ بحيث
      $R\sigma\sigma'$، إما~$\Delta(\sigma') \Proves/ !B$ أو
      $\Delta(\sigma') \Proves !C$. وبفرض الاستقراء يعني ذلك أنه كلما
      كان~$R\sigma\sigma'$، فإما
      $\mSat/{M(\Delta)}{!B}[\sigma']$ أو
      $\mSat{M(\Delta)}{!C}[\sigma']$؛ أي
      $\mSat{M(\Delta)}{\indfrm}[\sigma]$.}
  \end{enumerate}
\end{proof}

\end{document}
